\chapter{Glaucoma Surgery}
\section*{Overview}
Glaucoma surgery lowers intraocular pressure by improving fluid drainage or reducing fluid production to protect the optic nerve.
\section*{Common Instruments and Equipment}
The exact setup varies by technique and setting. Common items may include:
\begin{itemize}[nosep,leftmargin=*]
  \item operating microscope and ophthalmic microsurgical set
  \item procedure-specific implant or drainage device when used
  \item viscoelastic and irrigation-aspiration equipment
\end{itemize}

\section*{Why It Is Done}
It is considered when medicines or laser treatment do not adequately control pressure or when disease threatens vision.
\section*{How It Is Performed}
Options include trabeculectomy, glaucoma drainage-device implantation, minimally invasive glaucoma surgery, and cyclophotocoagulation. The choice depends on glaucoma type and severity.
\section*{Contraindications}
No single contraindication applies to every patient; the treating team assesses whether the procedure is suitable for the individual and the reason for treatment.

\section*{Precautions}
The care team reviews relevant health conditions, medicines, allergies, and procedure-specific risks, and explains any changes needed for the individual.

\section*{Risks and Recovery}
Bleeding, infection, pressure that is too low or high, scarring, cataract, inflammation, and vision loss are possible. Frequent postoperative examinations and drops are required.
\section*{References}
\begin{enumerate}\item National Eye Institute. Glaucoma Surgery. 2025.\item Current Medical Diagnosis and Treatment. McGraw-Hill; 2025.\end{enumerate}
\clearpage
